\section{Conclusión}
% Que puedo decir con base en lo que obtuve

En conclusión, la implementación de algoritmos genéticos para encontrar el mínimo de la 
función de Rosenbrock ha demostrado ser efectiva. Durante el desarrollo, se ha demostro
una convergencia hacia el mínimo global, lo que indica que los operadores de selección, cruza
y mutación están funcionando adecuadamente. Además, se ha observado que la diversidad
genética en la población juega un papel crucial en la exploración del espacio de búsqueda.
Los resultados obtenidos nos indican que, con un ajuste adecuado de los parámetros del algoritmo,
es posible mejorar aún más la eficiencia y precisión en la búsqueda del mínimo. Esto abre la
puerta a futuras investigaciones que busquen optimizar estos parámetros y explorar nuevas
variantes del algoritmo genético para abordar problemas similares.

Tambien cabe la pena mencionar que la función de Rosenbrock, es sencilla de implementar y
visualizar, lo que la convierte en un excelente caso de estudio para entender el comportamiento
en este caso podriamos haber explorado funciones mas complejas para ver como se comporta el algoritmo
en esos casos, pero para los fines de este reporte, la función de Rosenbrock ha sido suficiente.

Por otro lado la mutación ha demostrado ser un operador crucial para mantener la diversidad en la población,
evitando la convergencia prematura a óptimos locales. Sin embargo, es importante considerar que 
evitar la mutacion en este caso en especifico podria haber llevado a una convergencia más lenta
pero mas sencilla, lo cual puede ser un punto a considerar en futuros experimentos.

Es decir, es importante seguir explorando y ajustando los algoritmos genéticos para usar los parametros
correctos es decir cuando sean necesarios y cuando no, para maximizar su eficiencia en la búsqueda de soluciones óptimas.